\section{Experiments}

\subsection{Metrics}

We report three complementary metrics.
\emph{Corrupted accuracy} measures task performance when auxiliary text is misleading.
\emph{Incorrect aux-hit} measures the fraction of incorrect predictions that appear in the corrupted auxiliary text.
Lower incorrect aux-hit indicates less text-following among remaining errors, but it is not sufficient by itself: a model could reduce aux-hit by hallucinating unrelated answers.
Therefore, we interpret it together with corrupted accuracy and preservation on match/irrelevant conditions.
We also report corrected/regressed counts relative to the base model when available.

\subsection{DocVQA Main Results}

DocVQA is our main testbed.
We use Qwen3-VL-8B-Instruct as the primary model and evaluate on a held-out corrupted split.
Table~\ref{tab:docvqa-main} shows that text-following span DPO substantially improves both accuracy and text-bias behavior.
Increasing LoRA rank from 8 to 16 improves the result further, while the multi-condition error-mined variant gives the best overall performance.

\begin{table}[t]
\centering
\caption{DocVQA corrupted held-out results with Qwen3-VL-8B.}
\label{tab:docvqa-main}
\resizebox{\linewidth}{!}{
\begin{tabular}{lccc}
\toprule
Method & Acc. & Incorrect aux-hit & Corr./Reg. \\
\midrule
Base & 0.715 & 89.5\% & -- \\
TF-span DPO, r8 & 0.900 & 55.0\% & 37 / 0 \\
TF-span DPO, r16 & 0.915 & 47.1\% & 40 / 0 \\
Multi-cond. error DPO & \textbf{0.920} & \textbf{43.8\%} & 41 / 0 \\
\bottomrule
\end{tabular}
}
\end{table}

The multi-condition result suggests that the highest-performing variant is not corrupted-only.
However, the corrupted filter remains essential: corrupted rows still require the prediction to be copied from corrupted auxiliary text.
The additional match, irrelevant, and no-text errors act as weak condition-specific regularizers rather than replacing the text-following signal.

\subsection{DPO versus GRPO}

We compare DPO with GRPO-only and DPO-to-GRPO variants on the same disjoint DocVQA split.
GRPO samples multiple answers per prompt and optimizes a reward that encourages correctness and penalizes copying corrupted text.
Despite this denser reward, GRPO does not move the held-out greedy behavior meaningfully.
Table~\ref{tab:grpo} shows that DPO is the effective intervention in our setting.

\begin{table}[t]
\centering
\caption{GRPO-style optimization does not improve over DPO on DocVQA.}
\label{tab:grpo}
\resizebox{\linewidth}{!}{
\begin{tabular}{lccc}
\toprule
Variant & Init & Acc. & Incorrect aux-hit \\
\midrule
Base & -- & 0.715 & 89.5\% \\
GRPO-only, 100 steps & base & 0.710 & 89.7\% \\
GRPO-only, 300 steps & base & 0.705 & 89.8\% \\
GRPO-only, dense reward & base & 0.705 & 89.8\% \\
TF-span DPO & base & \textbf{0.900} & \textbf{55.0\%} \\
DPO $\rightarrow$ GRPO & DPO & \textbf{0.900} & 60.0\% \\
\bottomrule
\end{tabular}
}
\end{table}

DPO-to-GRPO changes only a few held-out predictions relative to DPO and does not improve accuracy.
This supports our hypothesis: when high-precision chosen/rejected pairs are available, direct pairwise optimization provides a clearer learning signal than free-form sampled-answer rewards.

\subsection{VQAv2 Transfer Check}

We repeat the disjoint split protocol on VQAv2.
The effect is positive but weaker than on DocVQA.
As shown in Table~\ref{tab:vqav2}, span DPO improves corrupted accuracy by 10 points and reduces incorrect aux-hit, but match and irrelevant preservation drop slightly.
Adding auxiliary SFT preservation does not recover this drop and slightly worsens corrupted performance.

\begin{table}[t]
\centering
\caption{VQAv2 transfer results with Qwen3-VL-8B.}
\label{tab:vqav2}
\resizebox{\linewidth}{!}{
\begin{tabular}{lcccc}
\toprule
Variant & Corr. Acc. & Aux-hit & Match & Irr. \\
\midrule
Base & 0.545 & 65.9\% & 0.92 & 0.77 \\
Span DPO & \textbf{0.645} & \textbf{40.8\%} & 0.90 & 0.75 \\
Span DPO + SFT & 0.635 & 42.5\% & 0.90 & 0.75 \\
\bottomrule
\end{tabular}
}
\end{table}

\subsection{Model Transfer}

We test whether the same mitigation extends beyond Qwen3-VL.
Table~\ref{tab:model-transfer} summarizes Qwen2-VL-7B and LLaVA-Next results on DocVQA.
Qwen2-VL benefits substantially, improving from 0.570 to 0.775 corrupted accuracy.
LLaVA-Next, however, has very low base DocVQA accuracy under our prompt/template, and DPO provides little benefit.

\begin{table}[t]
\centering
\caption{Model transfer on DocVQA. DPO helps Qwen2-VL but not LLaVA-Next under our setting.}
\label{tab:model-transfer}
\resizebox{\linewidth}{!}{
\begin{tabular}{llccc}
\toprule
Model & Setting & Acc. & Aux-hit & Corr./Reg. \\
\midrule
Qwen2-VL-7B & base & 0.570 & 93.0\% & -- \\
Qwen2-VL-7B & span DPO & \textbf{0.775} & 86.7\% & 41 / 0 \\
\midrule
LLaVA-N-7B & base & 0.105 & 90.5\% & -- \\
LLaVA-N-7B & span DPO & 0.100 & 91.7\% & 1 / 2 \\
\midrule
LLaVA-N-13B & base & 0.115 & 92.7\% & -- \\
LLaVA-N-13B & span DPO & 0.130 & 95.4\% & 3 / 0 \\
\bottomrule
\end{tabular}
}
\end{table}

These results show that the method is not simply a universal post-hoc fix.
It requires a base model that can already solve the task visually; DPO then teaches the model to resist misleading text.
